% =============================================================
\subsection{IoT — Justificativa dos Sensores para a Kodalabs}
% =============================================================

A adoção de dispositivos IoT na infraestrutura da Kodalabs não é uma escolha
meramente tecnológica, mas uma decisão estratégica diretamente vinculada ao perfil
operacional da empresa. Como consultoria de engenharia e ciência de dados, a
Kodalabs mantém equipamentos de alto valor computacional em operação contínua —
servidores, workstations de processamento intensivo e infraestrutura de rede —,
cujo funcionamento ininterrupto é condição essencial para a entrega dos serviços
aos clientes. Nesse contexto, os sensores IoT implantados cumprem três funções
críticas: \textbf{proteção patrimonial}, \textbf{continuidade operacional} e
\textbf{rastreabilidade de eventos físicos}.

\subsubsection{Sensor de Fumaça}

O sensor de fumaça é o dispositivo de maior criticidade para a Kodalabs. A empresa
opera com infraestrutura de servidores em regime 24/7, e ambientes de alta densidade
de equipamentos eletrônicos são particularmente suscetíveis a princípios de
incêndio causados por superaquecimento, falhas em fontes de alimentação ou curtos
em cabeamentos. Em uma consultoria de dados, um incêndio na sala de servidores
representa não apenas perda de hardware — substituível —, mas potencial destruição
de dados de clientes em processamento, o que pode resultar em violação de contratos,
penalidades por descumprimento da LGPD e dano irreparável à reputação da empresa.

O sensor de fumaça integrado à plataforma IoT do Data Center permite a detecção
precoce de anomalias, acionando alertas automáticos para a equipe de TI antes que
o evento se agrave. Isso reduz drasticamente o tempo de resposta em comparação a
sistemas de detecção convencionais não integrados à rede.

\subsubsection{Sensor de Movimento}

O sensor de movimento, implantado em São Carlos, atende a uma necessidade
específica da filial administrativa: o controle de acesso a áreas restritas fora
do horário comercial. A unidade de São Carlos concentra dados financeiros,
contratos com clientes e informações estratégicas da área comercial. O acesso
físico não autorizado a essa infraestrutura representa risco tanto de espionagem
corporativa quanto de sabotagem intencional.

A integração do sensor de movimento com a plataforma IoT central permite o registro
automático de eventos de presença, com \textit{timestamp} auditável, geração de
alertas em tempo real para os responsáveis de TI e eventual integração futura com
sistemas de controle de acesso físico (travas eletrônicas, câmeras IP). Para uma
empresa que trata dados sensíveis de clientes sob as obrigações da LGPD, a
rastreabilidade de quem acessou fisicamente os ambientes de infraestrutura é
também um requisito de conformidade.

\subsubsection{Atuadores de Ventilação}

Os atuadores de ventilação (fans) presentes em ambas as filiais respondem a um
desafio técnico concreto da Kodalabs: o gerenciamento térmico dos ambientes de
processamento. Workstations de engenharia de dados e ciência de dados executam
cargas computacionais intensas — pipelines ETL, treinamento de modelos de
\textit{machine learning}, processamento de grandes volumes de dados —, gerando
calor significativo de forma intermitente e imprevisível.

O controle manual de ventilação seria ineficiente e dependeria de presença física
constante. Com os atuadores integrados à plataforma IoT, é possível programar
regras automáticas (e.g., acionar ventilação adicional quando sensor de temperatura
exceder limite definido) ou controlar os dispositivos remotamente via interface web.
Isso reduz o risco de falha por superaquecimento, prolonga a vida útil dos
equipamentos e elimina a dependência de intervenção humana para eventos térmicos
que ocorrem fora do horário de expediente.

Em síntese, os três tipos de sensores escolhidos — fumaça, movimento e ventilação —
não foram selecionados aleatoriamente: cada um endereça um risco operacional real
e documentado do modelo de negócio da Kodalabs, tornando a infraestrutura IoT um
componente funcional da gestão de riscos da empresa.

% =============================================================
\subsection{VLANs — Justificativa da Segmentação para a Kodalabs}
% =============================================================

A divisão da rede em VLANs distintas para cada departamento reflete diretamente
a heterogeneidade funcional e de risco que caracteriza as equipes da Kodalabs.
Embora todos os colaboradores pertençam à mesma organização, os perfis de tráfego,
os níveis de confiabilidade dos dispositivos e as exigências de segurança variam
substancialmente entre departamentos. Colocar todas as equipes na mesma rede plana
seria equivalente a não ter segmentação alguma: um único dispositivo comprometido
teria visibilidade irrestrita sobre toda a infraestrutura.

\subsubsection{Por que separar Engenharia de Dados e Data Science?}

As equipes técnicas de Campinas (VLANs~10 e~20) geram volumes expressivos de
tráfego interno: transferências de datasets, sincronização de repositórios Git,
execução distribuída de pipelines Spark e comunicação entre nós de treinamento de
modelos. Se esse tráfego não fosse isolado, ele consumiria banda e aumentaria a
taxa de colisões de broadcast nas redes administrativas, degradando a experiência
de usuários que realizam tarefas leves como envio de e-mail e acesso a sistemas ERP.
A separação em VLANs distintas garante que o tráfego intensivo das equipes técnicas
fique confinado ao seu próprio domínio de broadcast, sem impacto nos demais
segmentos.

\subsubsection{Por que isolar a equipe de TI?}

A VLAN~30 de TI em Campinas e a VLAN~70 de TI em São Carlos concentram os
dispositivos com maior nível de privilégio na rede: estações de gerenciamento com
acesso SSH a roteadores e switches, ferramentas de monitoramento como SNMP e
\textit{syslog}, e credenciais administrativas de sistemas. Misturar esses
dispositivos com VLANs de usuários comuns criaria um vetor de ataque direto: um
colaborador mal-intencionado ou um dispositivo comprometido na VLAN de Engenharia
poderia, em uma rede plana, tentar acessar as ferramentas de gerenciamento da
equipe de TI por varredura simples de portas. O isolamento em VLAN dedicada, com
ACLs específicas, restringe esse acesso apenas aos dispositivos autorizados.

\subsubsection{Por que uma VLAN dedicada para servidores?}

A VLAN~40 de Campinas isola o servidor \texttt{SRV-CPN-DATA} em um segmento
de acesso controlado. Servidores que hospedam dados de clientes — especialmente
sob as obrigações da LGPD — não devem ser acessíveis diretamente por todos os
segmentos da rede. Ao concentrá-los em VLAN própria, é possível aplicar ACLs
granulares que determinam exatamente quais VLANs podem iniciar conexões aos
serviços (HTTP, FTP, SMTP) e em quais portas. Isso implementa o princípio de
\textit{least privilege} na camada de rede: nenhum dispositivo acessa o servidor
além do estritamente necessário para sua função.

\subsubsection{Por que isolar visitantes?}

As VLANs~50 (Campinas) e~90 (São Carlos) segregam dispositivos externos —
notebooks pessoais de clientes, smartphones de visitantes, equipamentos de
prestadores de serviço — que não passaram por nenhum processo de homologação de
segurança da Kodalabs. Esses dispositivos podem carregar malware, estar
desatualizados ou ser utilizados por pessoas com interesses contrários à empresa.
Em uma rede plana, um visitante conectado ao Wi-Fi teria acesso potencial a todos
os outros dispositivos da rede. Com a VLAN de visitantes, o acesso é restrito
exclusivamente à saída para a internet, sem visibilidade sobre qualquer recurso
corporativo interno.

\subsubsection{Separação geográfica: Campinas vs.\ São Carlos}

A adoção de blocos de endereçamento distintos — 192.168.10.0/24 para Campinas e
192.168.60.0/24 para São Carlos — não é apenas uma convenção de organização. Ela
permite que as tabelas de roteamento e as ACLs sejam escritas de forma
geograficamente semântica: uma regra que referencie \texttt{192.168.60.0/24}
aplica-se inequivocamente à filial de São Carlos, facilitando auditoria, diagnóstico
e expansão futura. Quando uma nova filial for adicionada (e.g., Ribeirão Preto com
bloco 192.168.80.0/24), a identidade geográfica de cada bloco se mantém clara sem
necessidade de reestruturação do esquema de endereçamento.

% =============================================================
\subsection{Segurança — Justificativa das ACLs para a Kodalabs}
% =============================================================

As Listas de Controle de Acesso (ACLs) implementadas no Projeto~2 são o mecanismo
que transforma a segmentação por VLANs de uma medida de organização em uma medida
efetiva de segurança. VLANs por si só separam domínios de broadcast na camada~2,
mas o tráfego entre VLANs é roteado na camada~3 pelo roteador. Sem ACLs, qualquer
dispositivo poderia potencialmente alcançar qualquer outro segmento da rede,
bastando que exista uma rota válida. As ACLs definem explicitamente quem pode falar
com quem, aplicando o princípio de \textit{least privilege} no nível do roteamento.

\subsubsection{ACL de Proteção da Infraestrutura IoT}

A ACL aplicada ao servidor \texttt{SRV-IOT} (192.168.100.3) é a medida de
segurança mais crítica do Projeto~2, e sua justificativa combina argumentos
técnicos e de segurança física.

Do ponto de vista técnico, os dispositivos IoT controlados pelo servidor — sensores
de fumaça, detectores de movimento e atuadores de ventilação — são
\textbf{atuadores físicos} além de simples coletores de dados. Acesso não
autorizado ao servidor IoT poderia permitir: desativação remota de sensores de
fumaça (eliminando proteção contra incêndio); geração de falsos alertas de movimento
(esgotando a atenção da equipe de segurança); ou manipulação dos atuadores de
ventilação (causando superaquecimento intencional de equipamentos). Essas ameaças
têm consequências físicas reais, não apenas digitais.

Do ponto de vista corporativo, a Kodalabs é uma empresa de consultoria de dados
que preza pela confiabilidade de sua infraestrutura perante clientes. Um incidente
causado por acesso indevido à plataforma IoT — especialmente se explorado por um
visitante conectado ao Wi-Fi — representaria falha grave de governança, com
potencial impacto reputacional e legal.

A política de ACL implementada é objetiva:

\begin{itemize}
    \item \textbf{Permitido:} sub-rede 192.168.10.128/27 (VLAN~30, TI de Campinas)
    — os responsáveis técnicos pela infraestrutura de Campinas têm acesso legítimo
    e necessário à plataforma IoT para configuração, monitoramento e resposta a
    incidentes.
    \item \textbf{Permitido:} sub-rede 192.168.60.64/26 (VLAN~70, TI de São Carlos)
    — a equipe técnica local de São Carlos precisa gerenciar os sensores de sua
    própria unidade e acompanhar o painel centralizado de monitoramento.
    \item \textbf{Negado explicitamente:} todas as demais sub-redes, incluindo
    VLANs de visitantes (192.168.10.168/29 e 192.168.60.160/28), VLANs
    administrativas, VLANs de engenharia e Data Science. Nenhum desses perfis
    tem justificativa funcional para interagir diretamente com a plataforma IoT.
\end{itemize}

A negação explícita dos visitantes merece destaque especial: visitantes conectados
ao Wi-Fi corporativo são, por definição, entidades não verificadas. Permitir que
um dispositivo de visitante acesse o servidor IoT — ainda que inadvertidamente —
seria equivalente a deixar a chave do painel de controle de segurança física em
local público. A ACL garante que, independentemente de qualquer falha de
configuração de VLAN, o servidor IoT só responda a endereços de origem pertencentes
às sub-redes de TI autorizadas.

\subsubsection{Alinhamento com a LGPD e Governança Corporativa}

As ACLs implementadas no Projeto~2 não são apenas medidas técnicas: elas constituem
controles de segurança documentáveis para fins de conformidade com a Lei Geral de
Proteção de Dados (LGPD). O art.~46 da LGPD exige que os agentes de tratamento
adotem medidas técnicas e administrativas aptas a proteger os dados pessoais de
acessos não autorizados. A segmentação por VLANs combinada com ACLs representa
exatamente esse conjunto de medidas: cada política de controle de acesso pode ser
auditada, documentada e apresentada como evidência de conformidade em eventuais
fiscalizações pela ANPD (Autoridade Nacional de Proteção de Dados).

% =============================================================
\section{Conclusão}
% =============================================================

O Projeto~2 da disciplina SSC0540 representa a maturação da infraestrutura de rede
da Kodalabs, evoluindo de uma arquitetura de filiais independentes para um ambiente
corporativo integrado, seguro e preparado para crescimento. As decisões técnicas
documentadas neste relatório — OSPF multi-área, segmentação por VLANs, controle de
acesso por ACLs e incorporação de IoT — não são fins em si mesmas, mas instrumentos
a serviço dos objetivos de negócio da empresa.

A adoção do OSPF multi-área garante que a comunicação entre Campinas, São Carlos e
o Data Center seja resiliente e escalável: a adição de novas filiais não exige
reconfiguração da topologia existente, e falhas de roteamento em uma área são
isoladas sem impactar as demais. O Data Center compartilhado, operando como
infraestrutura neutra na Área~0, centraliza serviços críticos sem favorecer
nenhuma das filiais, o que simplifica tanto a manutenção quanto a eventual expansão
para novas unidades.

A segmentação por VLANs traduz a diversidade funcional da Kodalabs em política de
rede: equipes técnicas com tráfego intensivo são isoladas de equipes administrativas,
servidores críticos operam em segmento de acesso controlado, e visitantes jamais
têm visibilidade sobre recursos internos. Esse modelo está diretamente alinhado com
os requisitos de conformidade da LGPD, que exige controles técnicos documentáveis
para proteção de dados pessoais tratados pela empresa.

A infraestrutura IoT introduzida no Projeto~2 não é um componente acessório, mas
uma extensão da gestão de riscos operacionais da Kodalabs. Sensores de fumaça,
detectores de movimento e atuadores de ventilação endereçam ameaças reais ao
ambiente físico onde opera a infraestrutura de dados da empresa. A proteção desses
dispositivos por ACLs específicas demonstra que segurança não é uma camada
adicionada ao final do projeto, mas um princípio que permeia cada decisão de
arquitetura desde sua concepção.

Como perspectivas de evolução futura, recomenda-se: a implementação de servidor
DHCP secundário com failover automático para aumentar a resiliência dos serviços;
a adoção de autenticação OSPF MD5 nos enlaces de backbone para proteção contra
injeção de rotas fraudulentas; a expansão da infraestrutura IoT com sensores de
temperatura integrados aos sistemas de ventilação para controle automatizado mais
preciso; e a implementação de \textit{syslog} centralizado no Data Center para
auditoria unificada de eventos de todas as áreas da rede.

A Kodalabs encerra o Projeto~2 com uma infraestrutura que não apenas suporta suas
operações atuais, mas está arquitetada para crescer com a empresa — mantendo
segurança, rastreabilidade e desempenho como propriedades estruturais, não como
adições pontuais.